\begin{textsample}{POS Dim 1 – 2021\_c4 – Score -4.00 – 2021\_c4/enade\_2021\_sist\_inf\_306a025769.txt}  \label{ex:f1_pos_372}
CAPÍTULO 4 PERCEPÇÃO DA PROVA Nas análises feitas neste capítulo, constam as percepções dos concluintes da Área de Sistemas de Informação ( Bacharelado ) sobre a prova aplicada no Enade/2021.
Essas percepções foram mensuradas por meio de nove questões que avaliaram desde o grau de dificuldade da prova até o tempo gasto para concluí -la.
As percepções sobre a prova foram relacionadas ao desempenho dos estudantes, à Grande Região de funcionamento do curso, à Categoria Administrativa e ao tipo de Organização Acadêmica da IES.
O questionário de percepção da prova encontra -se ao final do Anexo VIII, onde está a reprodução do exame.
O desempenho dos estudantes foi classificado em quatro quartos.
Para tanto, esse desempenho foi ordenado de forma ascendente.
O percentil 25, P25, também conhecido como primeiro quartil, é a nota de desempenho que deixa um quarto ( 25 \% ) dos valores observados abaixo e três quartos acima.
Na Figura 1, apresenta -se uma ilustração desse conceito.
O quarto inferior de desempenho é composto pelas notas abaixo do primeiro quartil.
Já o percentil 75, P75, também conhecido como terceiro quartil, é o valor para o qual há três quartos ( 75 \% ) dos dados abaixo e um quarto acima dele.
O quarto superior de desempenho é composto pelas notas iguais ou acima do terceiro quartil.
O percentil 50, P50, também conhecido como mediana, é o valor que divide as notas em dois conjuntos de igual tamanho.
O segundo quarto inclui os valores entre o primeiro quartil ( P25 ) e a mediana.
O terceiro quarto contém os valores entre a mediana ( P50 ) e o terceiro quartil ( P75 ).
Vale ressaltar que percentis, quartis e medianas são pontos que, não obrigatoriamente, pertencem ao conjunto original de dados, ao passo que os quartos são subconjuntos dos dados originais.
A seguir, serão apresentados gráficos com resultados selecionados, relativos às nove questões avaliadas por grupos de estudantes.
Nas barras dos gráficos, apresenta -se o percentual de alunos que assinalaram uma das opções ou a soma das porcentagens daqueles que assinalaram duas ( ou três ) delas.
Por exemplo, para as questões 1 e 2, nos gráficos, é apresentada a porcentagem total de participantes que assinalou as opções ( D ) Difícil e ( E ) Muito difícil.
Em cada barra, foram assinalados, também, os extremos do intervalo de confiança de 95 \% como linhas verticais unidas por uma linha horizontal na forma da letra H maiúscula.
O estimador de um parâmetro com um certo nível de confiança ( por exemplo, 95,0 \% ) deve conter o parâmetro no intervalo de confiança em 95 \% das vezes.
Na comparação entre os estimadores dos parâmetros de duas classes de uma dada categoria ( por exemplo : Norte e Nordeste nas Grandes Regiões, de primeiro e último quarto dentro de desempenho, de IES Públicas e Privada, ou de Universidades e Faculdades ), associados aos seus respectivos intervalos de confiança, diz -se que não há diferença estatisticamente significativa entre os parâmetros das duas categorias quando há interseção entre os intervalos de confiança, e que há diferença, se os intervalos de confiança são disjuntos ( para maiores informações vide Glossário ).
Nas Tabelas do Anexo II são apresentados os valores absolutos e a distribuição percentual18 das alternativas válidas das nove questões, segundo o mesmo recorte de desempenho dos alunos e da Grande Região de funcionamento do curso, de Categoria Administrativa e do tipo de Organização Acadêmica da IES. 4.1 GRAU DE DIFICULDADE DA PROVA 4.1.1 COMPONENTE DE FORMAÇÃO GERAL Ao avaliarem “ Qual o grau de dificuldade desta prova na parte de Formação Geral? ” ( Questão 1 ), 24,8 \% do grupo de inscritos e presentes optaram pelas alternativas Difícil ou Muito difícil.
Entretanto, para mais da metade dos estudantes ( 62,8 \% ), o Componente de Formação Geral da prova foi avaliado como tendo grau de dificuldade Médio ( Gráfico 4.1, Gráfico 4.2 e, no Anexo II, a Tabela II.1 ).
O percentual de estudantes que consideraram a prova Difícil ou Muito difícil foi maior na região Sul, onde a proporção foi de 29,3 \%, enquanto a de menor proporção foi a Sudeste, com 23,1 \%.
No Gráfico 4.1, é possível observar que as diferenças entre a região Sudeste e as regiões Sul e Centro-Oeste são estatisticamente significativas.
Nas Grandes Regiões, a proporção de presentes à prova que avaliaram o Componente de Formação Geral como de grau de dificuldade Médio esteve entre 58,5 \%, na região Sul, e 64,6 \%, na região Nordeste.
O percentual de alunos que consideraram o componente de Formação Geral da prova Difícil ou Muito difícil foi decrescente em função dos quartos de desempenho : 32,9 \%, no primeiro quarto, e 16,9 \%, no quarto, o grupo de melhor desempenho na prova.
Nos quartos de desempenho intermediários, a proporção de alunos que consideraram a prova Difícil ou Muito difícil foi de 27,6 \%, no segundo quarto, e de 21,9 \%, no terceiro quarto.
As diferenças entre todos os quartos de desempenho são estatisticamente significativas.
Para todos os quartos de desempenho, a alternativa modal para essa pergunta foi a Médio, com 58,0 \% e 66,1 \% dos respondentes nos quartos extremos, primeiro e quarto, respectivamente.
O percentual de estudantes que consideraram o componente de Formação Geral da prova Difícil ou Muito difícil foi maior entre os estudantes de CEFET/IF, alcançando 29,1 \%, enquanto entre os estudantes das Universidades, a Organização Acadêmica com a menor incidência, alcançou 22,8 \%.
No Gráfico 4.3, é possível observar que as diferenças de Universidades para Faculdades e CEFET/IF são estatisticamente significativas.
Nas Organizações Acadêmicas, a proporção de presentes à prova que avaliaram o Componente de Formação Geral como de grau de dificuldade Médio esteve entre 60,0 \%, nos CEFET/IF e Universidades, e 64,1 \% nas Universidades ( ver também a Tabela II.2, no Anexo II ).
O percentual de estudantes que consideraram o componente de Formação Geral da prova Difícil ou Muito difícil foi maior entre os estudantes de cursos de IES Privadas ( 25,6 \% ), sem diferença estatisticamente significativa da proporção registrada em IES Públicas ( 23,3 \% ).
Já a proporção de presentes à prova que avaliaram o Componente de Formação Geral como de grau de dificuldade Médio foi de 63,8 \%, entre os estudantes de IES Públicas, e de 62,3 \%, entre os de IES Privadas ( ver Gráfico 4.4 ). 4.1.2 COMPONENTE DE CONHECIMENTO ESPECÍFICO Ao responderem à Questão 2 – “ Qual o grau de dificuldade desta prova na parte de Componente Específico? ” – 41,6 \% do grupo de estudantes classificaram -na como Difícil ou Muito difícil.
Além disso, o Componente modal de Conhecimento Específico da prova foi considerado de grau de dificuldade Médio por 53,8 \% dos alunos ( Gráfico 4.5, Gráfico 4.6, e, no Anexo II, a Tabela II.3 ).
A análise das respostas dos estudantes quanto ao grau de dificuldade do Componente de Conhecimento Específico da prova, agregado por Grande Região, mostra diferenças estatisticamente significativas entre proporções de alunos que a avaliaram como Difícil ou Muito difícil.
Tais diferenças são observadas entre a região com a maior proporção, a Centro-Oeste ( 47,6 \% ), e as regiões Nordeste ( 41,1 \% ) e Sudeste ( 39,6 \% ), a com a menor proporção.
O percentual de alunos que classificaram o grau de dificuldade como Médio, no Componente de Conhecimento Específico, variou de 48,3 \% a 55,6 \%, nas regiões Centro-Oeste e Sudeste, respectivamente.
Considerando -se a avaliação da dificuldade das questões do Componente de Conhecimento Específico da prova de acordo com o desempenho dos estudantes, observa -se uma tendência decrescente dos que consideraram a parte específica Difícil ou Muito difícil, variando de 38,6 \% ( último quarto ) a 43,1 \% ( primeiro quarto ).
As demais proporções foram de 43,6 \% e 41,0 \%, respectivamente, no segundo e terceiro quartos.
Observa -se que há diferença estatisticamente significativa dos resultados entre o quarto superior de desempenho e o primeiro quarto.
Já os que responderam que o grau de dificuldade das questões do Componente de Conhecimento Específico da prova foi Médio atingiram percentagens que variaram de 50,5 \%, no primeiro quarto, a 57,4 \%, no último.
Considerando -se a avaliação da dificuldade das questões do Componente de Conhecimento Específico da prova, o percentual de estudantes que as consideraram Difícil ou Muito difícil foi maior entre os estudantes de CEFET/IF, atingindo 48,7 \% ; já entre os de Universidades, a Organização Acadêmica com a menor incidência, essa proporção foi de 39,6 \%.
No Gráfico 4.7, é possível observar que as diferenças de Faculdades ( 46,5 \% ) e de CEFET/IF para Universidades e Centros Universitários ( 40,6 \% ) são estatisticamente significativas.
Nas Organizações Acadêmicas, a proporção de presentes à prova que consideraram o Componente de Conhecimento Específico como tendo grau de dificuldade Médio esteve entre 48,0 \%, nos CEFET/IF, e 55,6 \%, nas Universidades ( ver também Tabela II.4, no Anexo II ).
O percentual de estudantes que consideraram o Componente Específico da prova Difícil ou Muito difícil foi o mesmo entre os estudantes de cursos de IES Privadas e os de IES Públicas, 41,6 \%, sem diferença estatisticamente significativa entre os percentuais.
Já a proporção de presentes à prova que avaliaram esse componente como sendo de grau de dificuldade Médio foi de 54,3 \%, entre os estudantes de IES Públicas, e de 53,5 \% entre os de IES Privadas ( ver Gráfico 4.8 e Tabela II.4, no Anexo II ). 4.2 EXTENSÃO DA PROVA EM RELAÇÃO AO TEMPO TOTAL Indagados quanto à extensão da prova, em relação ao tempo total oferecido para a sua resolução ( Questão 3 ), os estudantes apontaram, com maior incidência, a alternativa que considerava a extensão Adequada, em todas as agregações consideradas ( Gráfico 4.9, Gráfico 4.10, e, no Anexo II, a Tabela II.5 ).
O percentual de alunos que responderam ser a extensão da prova Adequada foi de 59,5 \%.
Já 33,3 \% dos inscritos presentes consideraram que a prova foi Longa ou Muito longa, e 7,2 \% a avaliaram como Curta ou Muito curta.
Entre as Grandes Regiões, os que avaliaram a prova como Longa ou Muito longa, em relação ao tempo total destinado à sua resolução, atingiram uma porcentagem que variou de 30,9 \%, na região Norte, até 34,1 \%, na região Nordeste.
Não há diferença estatisticamente significativa entre as Grande Regiões.
Considerando -se o desempenho dos alunos, nota -se ainda que 58,0 \% consideraram a extensão da prova Adequada, no quarto de desempenho inferior, e 60,0 \%, no de melhor desempenho ( quarto superior ).
Nos quartos intermediários, essa proporção foi de 60,7 \%, no segundo quarto, e de 59,3 \%, no terceiro.
No Gráfico 4.10, observa -se que as proporções de estudantes que consideraram a prova Longa ou Muito longa não apresentam tendência em função dos quartos de desempenho, sendo a maior proporção dessa resposta no último quarto de desempenho ( 34,1 \% ), e a menor, no segundo quarto ( 31,8 \% ).
Pode -se constatar que não há diferença estatisticamente significativa entre os quartos de desempenho.
Considerando -se a extensão da prova, o percentual de estudantes que avaliaram a prova como Longa ou Muito longa, em relação ao tempo total destinado à sua resolução, foi maior entre os estudantes de CEFET/IF, alcançando 34,9 \%, enquanto entre os estudantes de Centros Universitários, a Organização Acadêmica com a menor incidência, esse índice atingiu 31,8 \%.
Entre os de Faculdades, o índice alcançou 32,7 \%, e entre os de Universidades, 33,9 \%.
No Gráfico 4.11, é possível observar que não há diferença estatisticamente significativa entre os quatro tipos de Organização Acadêmica.
A proporção de presentes à prova que consideraram a sua extensão Adequada, em relação ao tempo total destinado à sua resolução, esteve entre 57,8 \%, nos CEFET/IF, e 60,2 \%, nas Faculdades e nos Centros Universitários ( ver também Tabela II.6, no Anexo II ).
O percentual de estudantes que consideraram a prova Longa ou Muito longa, em relação ao tempo total destinado à sua resolução, foi maior entre os estudantes de cursos de IES Públicas ( 34,6 \% ), não havendo diferença estatisticamente significativa da proporção em relação aos de IES Privadas ( 32,6 \% ).
Já a proporção de presentes à prova que consideraram a extensão Adequada foi de 58,2 \%, entre os estudantes de IES Públicas, e de 60,1 \%, entre os de IES Privadas ( ver Gráfico 4.12 e Tabela II.6, no Anexo II ). 4.3 COMPREENSÃO DOS ENUNCIADOS DAS QUESTÕES 4.3.1 COMPONENTE DE FORMAÇÃO GERAL Em relação aos enunciados das questões do Componente de Formação Geral ( Questão 4 ), as opiniões foram positivas, já que 74,8 \% dos alunos avaliados consideraram todos ou a maioria dos enunciados das questões “ claros e objetivos ” ( Gráfico 4.13, Gráfico 4.14 e, no Anexo II, a Tabela II.7 ).
Na análise regional, a percentagem de estudantes que avaliaram que todos ou a maioria dos enunciados das questões do Componente de Formação Geral estavam “ claros e objetivos ” variou de 73,8 \%, na região Sul, a 76,7 \%, na região Centro-Oeste.
A análise das percepções dos estudantes sobre a clareza e a objetividade dos enunciados permite afirmar que todos, ou a maioria dos enunciados de questões relativas ao Componente de Formação Geral, foram considerados “ claros e objetivos ” entre a maior parte dos respondentes ( maior ou igual a 73,8 \%, em todas as regiões, e maior ou igual a 65,3 \%, em todos os quartos de desempenho ).
Segundo o desempenho, observa -se que a proporção dos que emitiram essa opinião cresce, conforme o desempenho aumenta, com diferenças estatisticamente significativas entre todos os quartos de desempenho.
No quarto superior, a clareza e a objetividade de todos ou da maioria dos enunciados das questões foram percebidas por 81,7 \%.
A proporção dos presentes que consideraram que todos ou a maioria dos enunciados das questões do Componente de Formação Geral estavam “ claros e objetivos ” foi maior entre os estudantes de Centros Universitários, chegando a 76,7 \%, enquanto entre os de Faculdades, a Organização Acadêmica com a menor incidência, esse índice atingiu 71,6 \%.
Entre os de Universidades, o índice alcançou 74,4 \%, e entre os de CEFET/IF, 75,4 \%.
No Gráfico 4.15, é possível observar que as diferenças entre Centros Universitários e Faculdades são estatisticamente significativas ( ver também Tabela II.8, no Anexo II ).
O percentual de estudantes que consideraram que todos ou a maioria dos enunciados das questões do Componente de Formação Geral estavam “ claros e objetivos ” foi ligeiramente maior entre os estudantes de cursos de IES Públicas ( 74,8 \% ), e sem diferença estatisticamente significativa da proporção registrada em IES Privadas ( 74,7 \% ).
Um maior detalhamento pode ser encontrado no Gráfico 4.16 e na Tabela II.8, no Anexo II. 4.3.2 COMPONENTE DE CONHECIMENTO ESPECÍFICO Em relação aos enunciados das questões do Componente de Conhecimento Específico da prova, para 75,8 \% dos estudantes avaliados da Área de Sistemas de Informação ( Bacharelado ), a clareza e a objetividade ( Questão 5 ) estavam presentes em todas ou na maioria das questões ( Gráfico 4.17 a Gráfico 4.20, e no Anexo II, as Tabelas II.9 e II.10 ).
A maioria dos estudantes de todas as Grandes Regiões brasileiras considerou “ claros e objetivos ” todos ou a maioria dos enunciados das questões do Componente de Conhecimento Específico da prova, percentual sempre maior ou igual a 75,1 \%.
As diferenças entre as Grandes Regiões não são estatisticamente significativas.
A proporção de estudantes que consideraram os enunciados das questões “ claros e objetivos ” apresenta uma tendência crescente em relação ao aumento de desempenho : mais elevada no quarto superior ( 81,2 \% ), se comparada ao quarto inferior de desempenho ( 67,3 \% ).
A única diferença estatisticamente não significativa entre os quartos de desempenho é entre o terceiro quarto e o quarto superior.
A proporção dos presentes que consideraram que todos ou a maioria dos enunciados das questões do Componente de Conhecimento Específico estavam “ claros e objetivos ” foi maior entre os estudantes de Centros Universitários, atingindo 76,9 \%.
Entre os estudantes de Faculdades e de CEFET/IF, as Organizações Acadêmicas com a menor incidência, essa proporção atingiu 74,6 \%, sendo intermediária a proporção alcançada entre os de Universidades ( 75,7 \% ).
No Gráfico 4.19, é possível observar que as diferenças entre os tipos de Organização Acadêmica não são estatisticamente significativas ( ver também Tabela II.10, no Anexo II ).
O percentual de estudantes que consideraram que todos ou a maioria dos enunciados das questões do Componente de Conhecimento Específico estavam “ claros e objetivos ” foi ligeiramente maior entre os estudantes de cursos de IES Privadas ( 75,9 \% ), sem diferença estatisticamente significativa da proporção em IES Públicas ( 75,5 \% ).
Ver Gráfico 4.20 e Tabela II.10, no Anexo II, para maior detalhamento. 4.4 SUFICIÊNCIA DAS INFORMAÇÕES/INSTRUÇÕES FORNECIDAS Ao se avaliarem as informações/instruções fornecidas para a resolução das questões ( Questão 6 ), constatou -se que 85,0 \% dos respondentes da Área de Sistemas de Informação ( Bacharelado ) de todo o Brasil afirmaram que essas eram até excessivas ou suficientes em todas ou na maioria das questões ( Gráfico 4.21, Gráfico 4.22 e, no Anexo II, a Tabela II.11 ).
Quanto à distribuição de respondentes pelas Grandes Regiões, observa -se que a proporção de estudantes que consideraram as informações/instruções fornecidas até excessivas ou suficientes em todas ou na maioria das questões foi sempre superior ou igual a 82,7 \%, chegando a 86,4 \%, na região Sul.
Não existem diferenças estatisticamente significativas entre as Grandes Regiões.
Levando -se em conta o desempenho dos participantes, observa -se uma tendência crescente das proporções de participantes que avaliaram as informações/instruções como até excessivas ou suficientes em todas ou na maioria das questões, havendo diferenças estatisticamente significativas entre o primeiro, o segundo e o quarto superior, bem como entre o terceiro e o quarto superior de desempenho, conforme mostrado no Gráfico 4.22.
O percentual foi mais elevado no quarto superior ( 90,5 \% ), percentual superior à média nacional ( 85,0 \% ).
Já no quarto inferior, a suficiência das informações/instruções declarada como até excessiva, em todas ou na maioria das questões foi percebida por 79,1 \% dos respondentes.
Considerando -se as informações/instruções fornecidas para a resolução das questões ( Questão 6 ), a proporção que considerou que eram até excessivas ou suficientes em todas ou na maioria das questões foi maior entre os estudantes de Centros Universitários, totalizando 86,5 \%.
As Faculdades foram a Organização Acadêmica em que houve a menor proporção : 81,3 \%.
No Gráfico 4.23, é possível observar que as diferenças de Faculdades em relação a Universidades ( 85,5 \% ) e a Centros Universitários são estatisticamente significativas ( ver também Tabela II.12, no Anexo II ).
O percentual de estudantes que consideraram que as informações/instruções fornecidas para a resolução das questões eram até excessivas ou suficientes em todas ou na maioria das questões foi ligeiramente maior entre os estudantes de cursos de IES Privadas ( 85,1 \% ), não havendo diferença estatisticamente significativa da proporção registrada em IES Públicas ( 84,9 \% ) ( ver Gráfico 4.24 e Tabela II.12, no Anexo II ). 4.5 DIFICULDADE ENCONTRADA AO RESPONDER À PROVA Indagados sobre as dificuldades com as quais se depararam ao responder à prova ( Questão 7 ), 19,1 \% dos estudantes apontaram o Desconhecimento do conteúdo.
Para 40,1 \%, a Forma diferente de abordagem do conteúdo foi indicada como dificuldade.
Já a Falta de motivação para fazer a prova foi a dificuldade apontada por 21,6 \% dos respondentes.
Considerando -se todo o Brasil, 14,9 \% dos respondentes afirmaram que não tiveram qualquer tipo de dificuldade para responder à prova ( Tabelas II.13 e II.14, no Anexo II ).
Nos Gráficos 4.25 a 4.28, são apresentados os percentuais de estudantes que apontaram o Desconhecimento do conteúdo como “ dificuldade ao responder à prova ”.
Na análise por Grandes Regiões, o percentual de inscritos e presentes que apontaram o Desconhecimento do conteúdo como “ dificuldade ao responder à prova ” não superou 21,7 \%.
Os percentuais variaram de 17,8 \%, na região Sul, a 21,7 \%, na região Norte.
A Forma diferente de abordagem do conteúdo foi a escolha modal dos estudantes, com percentuais que variaram de 39,0 \% ( região Sudeste ) a 45,0 \% ( região Norte ).
O percentual de alunos que citaram a Falta de motivação como dificuldade variou de 13,2 \% ( região Norte ) a 23,8 \% ( região Sul ).
Os que declararam não ter qualquer tipo de dificuldade para responder à prova variaram de 13,0 \%, nas regiões Norte e Nordeste, a 15,9 \%, na região Sudeste.
Em relação aos quartos de desempenho, o Desconhecimento do conteúdo foi a opção escolhida por 17,9 \% dos estudantes do quarto inferior e por 19,9 \% do terceiro quarto, caracterizando uma tendência crescente.
Não são observadas diferenças estatisticamente significativas entre os quartos de desempenho.
A alternativa modal entre os alunos, quando agregados pelos quartos de desempenho, para a dificuldade encontrada, foi causada pela Forma diferente de abordagem do conteúdo : 37,8 \%, no primeiro quarto, e 42,3 \% do segundo quarto, escolheram essa alternativa.
Na análise por Organização Acadêmica, o percentual de inscritos e presentes que apontaram o Desconhecimento do conteúdo como “ dificuldade ao responder à prova ” não superou 20,8 \%.
Os percentuais variaram de 16,7 \%, nas Faculdades, a 20,8 \%, nas Universidades.
A Forma diferente de abordagem do conteúdo foi a escolha modal dos estudantes, com percentuais que variaram de 36,9 \% ( Universidades ) a 44,8 \% ( Faculdades ).
O percentual de alunos que citaram a Falta de motivação como dificuldade variou de 16,3 \% ( Centros Universitários ) a 24,9 \% ( Universidades ).
Os que declararam não ter qualquer tipo de dificuldade para responder à prova variaram de 11,0 \%, nos CEFET/IF, a 16,7 \%, nas Faculdades.
No Gráfico 4.27, é possível observar que não há diferenças estatisticamente significativas entre os quatro tipos de Organização Acadêmica ( ver também Tabela II.14, no Anexo II ).
Em relação às Categorias Administrativas, o Desconhecimento do conteúdo foi a opção escolhida por 20,2 \% dos estudantes, nas IES Públicas, e por 18,5 \%, nas IES Privadas.
Não são observadas diferenças estatisticamente significativas entre essas duas categorias.
A alternativa modal entre os alunos, quando agregados pelas Categorias Administrativas, para a dificuldade encontrada, foi causada pela Forma diferente de abordagem do conteúdo : 38,5 \%, entre os de IES Públicas, e 40,9 \%, entre os de IES Privadas ( ver Gráfico 4.28 e Tabela II.14, no Anexo II ). 4.6 CONTEÚDOS DAS QUESTÕES OBJETIVAS DA PROVA Ao se analisarem os conteúdos das questões objetivas da prova ( Questão 8 ), verifica -se que um percentual muito pequeno dos estudantes avaliados, apenas 3,9 \%, afirmou que não estudou ainda a maioria desses conteúdos ( Gráficos 4.29 a Gráfico 4.32, e as Tabelas II.15 e II.16, no Anexo II ).
A maioria ( 65,0 \% ) afirmou ter estudado e aprendido muitos ou todos os conteúdos.
Na análise por Grande Região, a proporção de respondentes que escolheram a opção não estudou ainda a maioria desses conteúdos foi pequena.
Observa -se que, nas regiões Norte ( 5,8 \% ), Nordeste ( 4,0 \% ) e Centro-Oeste ( 4,3 \% ), as proporções foram ligeiramente maiores do que a da média nacional ( 3,9 \% ).
Nas regiões Sul e Sudeste, a situação foi inversa, com proporções mais baixas, correspondendo a, respectivamente, 3,7 \% e 3,8 \%.
Não são observadas diferenças estatisticamente significativas entre as Grandes Regiões.
Em todas as regiões, a maioria dos presentes afirmou ter estudado e aprendido muitos ou todos os conteúdos, com proporções variando entre 57,3 \%, na região Norte, a 67,8 \%, na região Sul.
Considerando -se separadamente as opiniões de estudantes dos quatro quartos de desempenho, observa -se que, no quarto inferior, 7,8 \% ofereceram como resposta que não estudou ainda a maioria desses conteúdos, sendo 1,6 \% os do quarto superior com a mesma resposta.
A diferença entre os alunos que optaram por esse motivo de dificuldade no quarto inferior em relação aos demais quartos de desempenho é estatisticamente significativa, bem como a diferença entre os do segundo quarto e os do terceiro e quarto de desempenho superior.
Tendo -se em conta o quarto superior, 77,5 \% dos alunos afirmaram ter estudado e aprendido muitos ou todos os conteúdos.
No outro extremo, no primeiro quarto, 49,9 \% optaram pelas mesmas categorias.
As proporções são crescentes com o desempenho.
Na análise por Organização Acadêmica, a proporção de respondentes que escolheram a opção não estudou ainda a maioria desses conteúdos foi pequena.
Observa -se que a proporção nos CEFET/IF foi menor do que a da média nacional ( 3,9 \% ).
Nas Faculdades, Centros Universitários e Universidades a situação foi inversa : a proporção foi mais alta ( 4,0 \%, 4,0 \% e 4,2 \%, respectivamente ).
No Gráfico 4.31, é possível observar que as diferenças entre os tipos de Organização Acadêmica não são estatisticamente significativas.
Em todos os tipos de Organização Acadêmica, a maioria dos presentes afirmou ter estudado e aprendido muitos ou todos os conteúdos, com proporções variando entre 62,2 \%, nas Faculdades, a 66,7 \%, nos CEFET/IF ( ver também Tabela II.16, no Anexo II ).
Em relação às Categorias Administrativas, observa -se que, nas IES Públicas, 3,5 \% ofereceram como resposta que não estudou ainda a maioria desses conteúdos, sendo 4,2 \% nas IES Privadas com a mesma resposta.
A diferença entre os alunos que optaram por esse motivo de dificuldade nas duas categorias não é estatisticamente significativa.
Considerando -se as IES Públicas, 68,8 \% dos alunos afirmaram ter estudado e aprendido muitos ou todos os conteúdos, sendo de 62,9 \% a proporção dos que escolheram essa resposta nas IES Privadas ( ver Gráfico 4.32 e Tabela II.16, no Anexo II ). 4.7 TEMPO GASTO PARA CONCLUIR A PROVA Ao responderem sobre o tempo de conclusão da prova ( Questão 9 ), mais de quatro quintos dos estudantes ( 70,0 \% ) afirmaram ter gasto Entre duas e quatro horas ( Gráfico 4.33 a Gráfico 4.36 e, no Anexo II, as Tabelas II.17 e II.18 ).
Considerando -se as cinco Grandes Regiões brasileiras, os percentuais dos que utilizaram Entre duas e quatro horas para finalizar a prova nas regiões Sudeste ( 68,8 \% ) e Sul ( 66,8 \% ) foram inferiores ao percentual nacional.
Nas demais Grandes Regiões, os percentuais de alunos que dispensaram Entre duas e quatro horas para concluir a prova foram mais elevados que a média : 75,1 \% na região Norte ; 73,9 \%, na região Nordeste ; e 71,1 \%, na região Centro-Oeste, como mostrado no Gráfico 4.33.
Há diferenças estatisticamente significativas das duas regiões com percentuais mais baixos – Sudeste e Sul – em relação às duas regiões com percentuais mais elevados – Norte e Nordeste.
Uma vez analisadas as alternativas escolhidas pelos estudantes que se situam nos diferentes quartos de desempenho, observa -se uma tendência crescente : uma proporção maior de participantes no quarto superior declarou ter gasto Entre duas e quatro horas para concluir a prova quando comparada com a proporção de participantes nos quartos inferiores.
Há diferenças estatisticamente significativas entre todos os quartos de desempenho.
Na análise por Organização Acadêmica, os percentuais dos que utilizaram Entre duas e quatro horas para finalizar a prova nos Centros Universitários ( 73,3 \% ), Faculdades ( 75,6 \% ) e CEFET/IF ( 74,9 \% ) foram superiores ao percentual nacional ( 70,0 \% ).
Nas Universidades, tal percentual foi inferior à média, 66,3 \%, e, como mostrado no Gráfico 4.35, há diferença estatisticamente significativa entre as Universidades e os outros três tipos de Organização Acadêmica.
Em relação às Categorias Administrativas, observa -se que, nas IES Públicas, 71,2 \% declararam ter gasto Entre duas e quatro horas para concluir a prova, sendo 69,4 \% nas IES Privadas com a mesma resposta.
A diferença entre os alunos que optaram por esse motivo de dificuldade nas duas categorias não é estatisticamente significativa ( ver Gráfico 4.36 e Tabela II.18, no Anexo II ).

% matched lemmas: 
\end{textsample}
